\chapter*{Resumen}

Este Trabajo Fin de Grado presenta OmniLitter, una plataforma multiplataforma
orientada a la gestión, recogida y análisis de datos sobre contaminación por
residuos en entornos terrestres y acuáticos. El proyecto surge de la necesidad de
digitalizar los procesos de muestreo y unificar la información obtenida en campo
para que pueda consultarse, trazarse y reutilizarse con criterios científicos.

La solución propuesta integra una aplicación móvil, un portal web y una API
centralizada dentro de una arquitectura común. Esta estructura permite cubrir el
ciclo completo de trabajo, desde la organización de campañas y la captura de
datos en campo hasta la consulta, explotación y exportación de la información
recopilada. El sistema se apoya en una taxonomía común de residuos, en el
registro geolocalizado de observaciones y en mecanismos de sincronización que
permiten trabajar en contextos con conectividad limitada.

A lo largo del trabajo se abordan el análisis de requisitos, el diseño de la
arquitectura y del modelo de datos, la planificación del desarrollo, la
implementación de los módulos principales, el despliegue y la validación mediante
pruebas. El resultado es un prototipo funcional que demuestra la viabilidad de
una plataforma unificada para apoyar tareas de investigación, seguimiento
ambiental y evaluación de medidas de reducción de residuos. OmniLitter
constituye así una base tecnológica extensible sobre la que pueden incorporarse
nuevas capacidades en trabajos futuros.

\vspace{.5cm}

\textbf{Palabras clave:} OmniLitter, monitorización de residuos, aplicación
móvil, portal web, geolocalización, trabajo sin conexión, sincronización,
análisis de datos ambientales.
